\chapter{Expected Utility}\label{ch:b2:expected-utility}

Offer a group of people a choice between \num{2400} for certain and a gamble that pays \num{2500} with
probability $0.33$, \num{2400} with probability $0.66$ and nothing with probability $0.01$. Most take
the sure \num{2400}. Now offer the same people a second choice: \num{2500} with probability $0.33$, or
\num{2400} with probability $0.34$. Most take the \num{2500}. Each choice looks sensible on its own.
Together they break the theory of choice under risk that economics uses everywhere else, from
insurance to pensions to the pricing of shares, and this chapter shows why. The chart below gives
the answers of the respondents of Kahneman and Tversky (1979), who asked the two questions with
prizes in Israeli pounds (the median family's net monthly income was then about \num{3000} pounds),
% ledger: b2:kt1979-income
and two more questions built the same way.

\begin{openingchart}{The Allais pairs. Percentage of respondents choosing the less risky option in each
  of two pairs of questions: first with a sure prize, then with both options scaled down. Expected
  utility says that each respondent makes the same choice twice, so the two bars of a pair should
  have the same height.}{Kahneman and Tversky (1979), Problems 1--4; 72 respondents for the first pair, 95 for the second}
\begin{tikzpicture}
\begin{axis}[omrecord, ybar, bar width=16pt, ymin=0, ymax=100, xmin=0.5, xmax=2.5,
    xtick={1,2}, xticklabels={{Problems 1 and 2}, {Problems 3 and 4}},
    ylabel={Percentage of respondents}, legend pos=north east, grid=none, ymajorgrids,
    nodes near coords, every node near coord/.append style={font=\footnotesize}]
  \addplot[fill=omDef!70, draw=omDef] table[x=pos, y=certain]
    {figdata/bachelor-2/01-expected-utility/out/allais.csv};
  \addplot[fill=omThm!45, draw=omThm] table[x=pos, y=scaled]
    {figdata/bachelor-2/01-expected-utility/out/allais.csv};
  \legend{With the sure prize, Scaled down}
\end{axis}
\end{tikzpicture}
\end{openingchart}

In the first pair, 82 per cent take the sure \num{2400}, but only 17 per cent take the matching
% figdata: allais (certain 82, scaled 17); ledger: b2:kt1979-p1, b2:kt1979-p2
option once the two prospects are scaled down. In the second pair the figures are 80 and 35 per cent.
% figdata: allais (certain 80, scaled 35); ledger: b2:kt1979-p3, b2:kt1979-p4
The gap is the size of the failure: at least 65 per cent of the respondents in the first pair, and
% figdata: allais (gap 65, 45)
45 per cent in the second, made a pair of choices that no expected-utility maximiser makes. The
chapter builds the theory first, as a theorem with four axioms, then measures attitudes to risk with
it, and ends with the places where it is tested and where it breaks.

\begin{recall}
From the school volumes: a gamble is described by its prizes and their probabilities; its expected
value is the probability-weighted average of the prizes; a person who prefers the expected value for
sure to the gamble itself is risk-averse; the sure amount she values as much as the gamble is its
\omterm{def:b2:expected-utility:certainty-equivalent}{certainty equivalent}. From mathematics: expectation, variance, concave functions and Jensen's
inequality, Taylor's formula to second order.
\end{recall}

\begin{notation}[Symbols of this chapter]\label{not:b2:expected-utility:symbols}
Amounts of money are $z$, prizes $z_1 < z_2 < \dots < z_n$. A \omterm{def:b2:expected-utility:lottery}{lottery} is $\ell = (\ell_1, \dots,
\ell_n)$, and $\ell'$, $\ell''$ are others; $\theta \in [0,1]$ is a mixing weight. The \omterm{def:b2:expected-utility:risk-premium}{risk premium} is
$\rho$, a constant of \omterm{def:b2:expected-utility:risk-aversion}{risk aversion} $a$ or $\gamma$, and the Arrow--Pratt coefficients are
$\mathcal{A}(z)$ and $\mathcal{R}(z)$. A claim comes with probability $\mu$; an amount of
insurance cover is $\chi$, bought at $\psi$ per unit of cover. (In this series $p$ is a price, $w$ a wage, $\pi$ inflation,
$\alpha$ a capital share, $P$ the price level and $C$ consumption, so none of them is used here.)
\end{notation}

\section{Lotteries}\label{sec:b2:expected-utility:lotteries}

The objects of choice are not goods but probability distributions over outcomes. Fix a finite set of
prizes $Z = \{z_1, \dots, z_n\}$: amounts of money, or any outcomes at all.

\begin{definition}[Lottery]\label{def:b2:expected-utility:lottery}
A \emph{lottery}\index{lottery} on the prizes $z_1, \dots, z_n$ is a vector $\ell = (\ell_1, \dots,
\ell_n)$ with $\ell_i \ge 0$ and $\sum_i \ell_i = 1$, where $\ell_i$ is the probability of winning
$z_i$. The set of lotteries is the simplex $\Delta = \{\ell \in \mathbb{R}^n : \ell_i \ge 0, \sum_i
\ell_i = 1\}$. The \emph{degenerate lottery} $\delta_i$ gives $z_i$ for certain. When the prizes are
amounts of money, $\mathbb{E}[\ell] = \sum_i \ell_i z_i$ is the lottery's expected value.
\end{definition}

\begin{definition}[Compound lottery]\label{def:b2:expected-utility:compound-lottery}
A \emph{compound lottery}\index{compound lottery} gives the \omterm{def:b2:expected-utility:lottery}{lottery} $\ell^{(k)}$ with probability
$\theta_k$, $k = 1, \dots, K$ ($\theta_k \ge 0$, $\sum_k \theta_k = 1$): first one draw picks a \omterm{def:b2:expected-utility:lottery}{lottery},
then a second draw picks the prize. Its \emph{reduced \omterm{def:b2:expected-utility:lottery}{lottery}} is $\sum_k \theta_k \ell^{(k)}$, which
gives $z_i$ with the total probability $\sum_k \theta_k \ell^{(k)}_i$. The mixture $\theta\ell +
(1-\theta)\ell'$ is the reduced \omterm{def:b2:expected-utility:lottery}{lottery} of the compound \omterm{def:b2:expected-utility:lottery}{lottery} that gives $\ell$ with probability
$\theta$ and $\ell'$ otherwise.
\end{definition}

Throughout, a person cares only about the reduced \omterm{def:b2:expected-utility:lottery}{lottery}: two compound lotteries with the same
reduced \omterm{def:b2:expected-utility:lottery}{lottery} are the same object of choice. This is an assumption, not a theorem, and it is part
of what the evidence tests.

\begin{example}[The Allais lotteries]\label{ex:b2:expected-utility:allais-lotteries}
Take the prizes $z_1 = 0$, $z_2 = \num{2400}$, $z_3 = \num{2500}$. The four options of the opening
chart's first pair are
\[
\ell^{A} = (0.01, 0.66, 0.33),\quad \ell^{B} = (0, 1, 0),\quad
\ell^{C} = (0.67, 0, 0.33),\quad \ell^{D} = (0.66, 0.34, 0).
\]
Let $\ell^{*} = (\tfrac{1}{34}, 0, \tfrac{33}{34})$, which pays \num{2500} with probability $33/34$.
Then
\[
\ell^{A} = 0.34\,\ell^{*} + 0.66\,\delta_2,\quad \ell^{B} = 0.34\,\delta_2 + 0.66\,\delta_2,\quad
\ell^{C} = 0.34\,\ell^{*} + 0.66\,\delta_1,\quad \ell^{D} = 0.34\,\delta_2 + 0.66\,\delta_1 .
\]
Check the first: $0.34 \times \tfrac{1}{34} = 0.01$ on $0$, $0.66$ on \num{2400}, $0.34 \times
\tfrac{33}{34} = 0.33$ on \num{2500}. Each pair mixes the same two lotteries, $\ell^{*}$ and the sure
\num{2400}, with a common consequence: \num{2400} with probability $0.66$ in the first pair, nothing
with probability $0.66$ in the second.
\end{example}

With three prizes a \omterm{def:b2:expected-utility:lottery}{lottery} is fixed by two numbers, $\ell_1$ and $\ell_3$, since $\ell_2 = 1 - \ell_1
- \ell_3$. The \emph{probability triangle} plots $\ell_1$ (the chance of the worst prize) across and
$\ell_3$ (the chance of the best) up. Its corners are the three sure prizes: $\delta_2$ at the origin,
$\delta_1$ at $(1, 0)$, $\delta_3$ at $(0, 1)$. Moving up or to the left shifts probability to better
prizes. The figure after \cref{prop:b2:expected-utility:parallel} draws the Allais lotteries in it.

\section{The expected-utility theorem}\label{sec:b2:expected-utility:theorem}

A person has a preference relation $\succsim$ over $\Delta$: $\ell \succsim \ell'$ reads ``$\ell$ is
at least as good as $\ell'$''; $\succ$ and $\sim$ are its strict part and its indifference, as for
bundles of goods in the Year 1 volume. Four axioms are imposed.

\begin{enumerate}
  \item \textbf{Completeness:} for all $\ell, \ell'$, $\ell \succsim \ell'$ or $\ell' \succsim \ell$.
  \item \textbf{Transitivity:} $\ell \succsim \ell'$ and $\ell' \succsim \ell''$ imply $\ell \succsim
        \ell''$.
  \item \textbf{Continuity:} for all $\ell, \ell', \ell''$ the sets $\{\theta \in [0,1] : \theta\ell +
        (1-\theta)\ell' \succsim \ell''\}$ and $\{\theta \in [0,1] : \ell'' \succsim \theta\ell +
        (1-\theta)\ell'\}$ are closed.
  \item \textbf{Independence}, defined next.
\end{enumerate}

\begin{definition}[Independence axiom]\label{def:b2:expected-utility:independence-axiom}
The preference $\succsim$ satisfies the \emph{independence axiom}\index{independence axiom} if, for
all lotteries $\ell, \ell', \ell''$ and every $\theta \in (0, 1]$,
\[
\ell \succsim \ell' \iff \theta\ell + (1-\theta)\ell'' \succsim \theta\ell' + (1-\theta)\ell'' .
\]
Mixing both lotteries with a third one, in the same proportion, leaves the ranking unchanged.
\end{definition}

Because the axiom is an equivalence, it holds for $\succ$ and for $\sim$ as well. Its logic is that of
the \omterm{def:b2:expected-utility:compound-lottery}{compound lottery}: with probability $1-\theta$ the person gets $\ell''$ whatever she chose, so only
the branch where $\ell$ and $\ell'$ differ should matter. Continuity rules out lexicographic tastes, such
as preferring any \omterm{def:b2:expected-utility:lottery}{lottery} with a smaller chance of death, however small the difference, whatever else
it offers.

\begin{definition}[Expected-utility representation]\label{def:b2:expected-utility:eu-representation}
A preference $\succsim$ on $\Delta$ has an \emph{expected-utility
representation}\index{expected-utility representation} if there is a function $u : Z \to \mathbb{R}$,
the \emph{utility of money} (or Bernoulli utility function), such that
\[
\ell \succsim \ell' \iff U(\ell) \ge U(\ell'), \qquad U(\ell) = \sum_{i=1}^{n} \ell_i\,u(z_i) .
\]
$U$ is the \emph{expected utility} of the \omterm{def:b2:expected-utility:lottery}{lottery}. For a random amount $\tilde z$ we write $U =
\mathbb{E}[u(\tilde z)]$.
\end{definition}

$U$ is linear in the probabilities: $U(\theta\ell + (1-\theta)\ell') = \theta U(\ell) + (1-\theta)
U(\ell')$. That one property is the whole content of the theory, and the theorem says that the four
axioms are exactly what it takes.

\begin{theorem}[Expected-utility theorem (von Neumann--Morgenstern)]\label{thm:b2:expected-utility:vnm}
\index{expected-utility theorem}
A preference relation on the lotteries over finitely many prizes is complete, transitive, continuous
and satisfies independence if and only if it has an \omterm{def:b2:expected-utility:eu-representation}{expected-utility representation}.
\end{theorem}

\begin{proof}
\emph{If.} Let $U(\ell) = \sum_i \ell_i u(z_i)$ represent $\succsim$. Completeness and transitivity
hold because $\ge$ on numbers has them. Since $U(\theta\ell + (1-\theta)\ell')$ is an affine function
of $\theta$, the sets of continuity are of the form $\{\theta : \text{an affine function of }\theta
\ge \text{a constant}\}$, which are closed intervals. For independence,
$U(\theta\ell + (1-\theta)\ell'') - U(\theta\ell' + (1-\theta)\ell'') = \theta\,(U(\ell) - U(\ell'))$,
which has the sign of $U(\ell) - U(\ell')$ when $\theta > 0$.

\emph{Only if.} Six steps.

\emph{Step 1: a best and a worst \omterm{def:b2:expected-utility:lottery}{lottery}.} The $n$ sure prizes are ranked by completeness and
transitivity; let $\bar\delta$ be a best and $\underline\delta$ a worst of them. We show $\bar\delta
\succsim \ell \succsim \underline\delta$ for every $\ell$, by induction on the number of prizes with
positive probability. One prize: $\ell$ is a sure prize. Otherwise, with $\ell_1 \in (0,1)$ say, write
$\ell = \ell_1 \delta_1 + (1-\ell_1)\ell^{(1)}$, where $\ell^{(1)}$ is $\ell$ conditional on not
winning $z_1$ and has fewer prizes, so $\bar\delta \succsim \ell^{(1)}$. Independence applied twice
gives
\[
\ell = \ell_1\delta_1 + (1-\ell_1)\ell^{(1)} \precsim \ell_1\bar\delta + (1-\ell_1)\ell^{(1)} \precsim
\ell_1\bar\delta + (1-\ell_1)\bar\delta = \bar\delta ,
\]
the first step with common part $\ell^{(1)}$, the second with common part $\bar\delta$. The other
inequality is the same argument. If $\bar\delta \sim \underline\delta$, every \omterm{def:b2:expected-utility:lottery}{lottery} is indifferent
to every other and a constant $u$ represents $\succsim$; assume from now on $\bar\delta \succ
\underline\delta$.

\emph{Step 2: more weight on the best is better.} Write $m(\theta) = \theta\bar\delta +
(1-\theta)\underline\delta$. If $1 \ge \theta > \theta' \ge 0$, then $m(\theta) \succ m(\theta')$.
Indeed, with $\beta = (\theta - \theta')/(1 - \theta') \in (0, 1]$, $m(\theta) = \beta\bar\delta +
(1-\beta) m(\theta')$. Independence (strict form, common part $\bar\delta$) gives $\bar\delta =
\theta'\bar\delta + (1-\theta')\bar\delta \succ m(\theta')$, and then (common part $m(\theta')$)
$m(\theta) = \beta\bar\delta + (1-\beta)m(\theta') \succ \beta m(\theta') + (1-\beta) m(\theta') =
m(\theta')$.

\emph{Step 3: every \omterm{def:b2:expected-utility:lottery}{lottery} has a calibrating weight.} Fix $\ell$. The sets $S^{+} = \{\theta :
m(\theta) \succsim \ell\}$ and $S^{-} = \{\theta : \ell \succsim m(\theta)\}$ are closed by
continuity, $1 \in S^{+}$ and $0 \in S^{-}$ by step 1, and $S^{+} \cup S^{-} = [0,1]$ by
completeness. An interval is connected, so two non-empty closed sets that cover it meet: there is
$\theta_\ell$ with $\ell \sim m(\theta_\ell)$. It is unique by step 2.

\emph{Step 4: the weight represents the preference.} Let $U(\ell) = \theta_\ell$. By transitivity,
$\ell \succsim \ell' \iff m(\theta_\ell) \succsim m(\theta_{\ell'})$, and by step 2 and completeness
this holds if and only if $\theta_\ell \ge \theta_{\ell'}$.

\emph{Step 5: $U$ is linear.} Since $\ell \sim m(\theta_\ell)$ and $\ell' \sim m(\theta_{\ell'})$,
independence in its indifference form, used once for each \omterm{def:b2:expected-utility:lottery}{lottery}, gives
\[
\theta\ell + (1-\theta)\ell' \sim \theta\, m(\theta_\ell) + (1-\theta)\ell' \sim \theta\, m(\theta_\ell)
+ (1-\theta)\, m(\theta_{\ell'}) = m\bigl(\theta\theta_\ell + (1-\theta)\theta_{\ell'}\bigr).
\]
By the uniqueness in step 3, $U(\theta\ell + (1-\theta)\ell') = \theta U(\ell) + (1-\theta)U(\ell')$.

\emph{Step 6: the expected-utility form.} Every \omterm{def:b2:expected-utility:lottery}{lottery} is a mixture of sure prizes, $\ell = \sum_i
\ell_i\delta_i$. Applying step 5 repeatedly (induction on $n$) gives $U(\ell) = \sum_i \ell_i
U(\delta_i)$. Set $u(z_i) = U(\delta_i)$.
\end{proof}

The proof also shows how to measure $u$: $u(z_i)$ is the probability $\theta$ of the best prize at
which the person is indifferent between $z_i$ for sure and the \omterm{def:b2:expected-utility:lottery}{lottery} $m(\theta)$. A utility of
money is a schedule of such answers.

\begin{proposition}[Uniqueness up to a positive affine map]\label{prop:b2:expected-utility:affine}
If $u$ gives an \omterm{def:b2:expected-utility:eu-representation}{expected-utility representation} of $\succsim$, then $v$ gives one too if and only if
$v = \kappa_0 + \kappa_1 u$ for some numbers $\kappa_0$ and $\kappa_1 > 0$.
\end{proposition}

\begin{proof}
If $v = \kappa_0 + \kappa_1 u$, then $\sum_i \ell_i v(z_i) = \kappa_0 + \kappa_1 \sum_i \ell_i u(z_i)$,
an increasing function of $U(\ell)$, so $v$ ranks lotteries as $u$ does. Conversely, suppose both
represent $\succsim$ and $\bar\delta \succ \underline\delta$ (otherwise both are constant and
$\kappa_1 = 1$, $\kappa_0 = v - u$ will do). Write $\bar u, \underline u, \bar v, \underline v$ for the
utilities of the best and worst prizes. For each prize, $\delta_i \sim m(\theta_i)$ for a unique
$\theta_i$, and both representations must give the two lotteries equal value:
\[
u(z_i) = \theta_i\bar u + (1-\theta_i)\underline u, \qquad v(z_i) = \theta_i\bar v +
(1-\theta_i)\underline v .
\]
Eliminate $\theta_i = (u(z_i) - \underline u)/(\bar u - \underline u)$: $v(z_i) = \underline v +
\kappa_1 (u(z_i) - \underline u)$ with $\kappa_1 = (\bar v - \underline v)/(\bar u - \underline u) >
0$.
\end{proof}

So the units and the zero of $u$ mean nothing; ratios of utility differences mean something. We may
always set $u$ of the worst prize to $0$ and $u$ of the best to $1$. A statement such as ``she is
twice as happy with \num{2500} as with \num{1000}'' has no content in this theory; ``the step from
\num{1000} to \num{2500} is worth three times the step from $0$ to \num{1000}'' has.

\begin{proposition}[Straight parallel indifference lines]\label{prop:b2:expected-utility:parallel}
Under expected utility with three prizes $z_1 < z_2 < z_3$ and $u(z_1) < u(z_2) < u(z_3)$, the
indifference curves in the probability triangle are parallel straight lines of slope
\[
\frac{\dd\ell_3}{\dd\ell_1} = \frac{u(z_2) - u(z_1)}{u(z_3) - u(z_2)} > 0,
\]
and utility rises towards the upper left.
\end{proposition}

\begin{proof}
With $\ell_2 = 1 - \ell_1 - \ell_3$, $U = u(z_2) - \ell_1\,(u(z_2) - u(z_1)) + \ell_3\,(u(z_3) -
u(z_2))$, an affine function of $(\ell_1, \ell_3)$. Its level sets are parallel lines, and setting
$\dd U = 0$ gives the slope. $U$ falls with $\ell_1$ and rises with $\ell_3$.
\end{proof}

\begin{omfigure}
\begin{tikzpicture}[econ, scale=6.2]
  \draw[econ axes, <->] (0,1.08) node[left]{$\ell_3$} -- (0,0) -- (1.08,0) node[below]{$\ell_1$};
  \draw[thick] (1,0) -- (0,1);
  \node[below left] at (0,0) {$\delta_2$};
  \node[below] at (1,0) {$\delta_1$};
  \node[left] at (0,1) {$\delta_3$};
  \foreach \c in {0.08,0.36,0.64,0.92} {
    \draw[guide] ({max(0,-\c/2)},{\c+2*max(0,-\c/2)}) -- ({(1-\c)/3},{\c+2*(1-\c)/3});
  }
  \draw[guide] (0.16,0) -- ({0.16+0.28},{0.56});
  \draw[guide] (0.44,0) -- ({0.44+0.187},{0.373});
  \node[eqpoint, label={[font=\footnotesize]above right:$\ell^{B}$}] (B) at (0,0) {};
  \node[eqpoint, label={[font=\footnotesize]left:$\ell^{A}$}] (A) at (0.01,0.33) {};
  \node[eqpoint, label={[font=\footnotesize]below:$\ell^{D}$}] (D) at (0.66,0) {};
  \node[eqpoint, label={[font=\footnotesize]right:$\ell^{C}$}] (C) at (0.67,0.33) {};
  \draw[omExo, very thick] (B) -- (A);
  \draw[omExo, very thick] (D) -- (C);
  \node[font=\footnotesize, align=left, anchor=west] at (0.52,0.78)
    {dotted: indifference lines\\of one utility of money,\\utility rising to the upper left};
\end{tikzpicture}

{\small The probability triangle, with $z_1 = 0$, $z_2 = \num{2400}$, $z_3 = \num{2500}$. The segments
from $\ell^{B}$ to $\ell^{A}$ and from $\ell^{D}$ to $\ell^{C}$ are parallel (both are the vector
$(0.01, 0.33)$). With straight parallel indifference lines, $\ell^{A}$ lies on the better side of the
line through $\ell^{B}$ exactly when $\ell^{C}$ lies on the better side of the line through
$\ell^{D}$: the two choices must go the same way.}
\end{omfigure}

\section{Attitudes to risk}\label{sec:b2:expected-utility:attitudes}

From here on the prizes are amounts of money and $u$ is increasing and continuous. The theorem extends
to lotteries with infinitely many outcomes (distributions of a random amount $\tilde z$ on an
interval) under a further continuity condition, which we admit: $U = \mathbb{E}[u(\tilde z)]$.

\begin{definition}[Risk aversion]\label{def:b2:expected-utility:risk-aversion}
A person with utility of money $u$ shows \emph{risk aversion}\index{risk aversion} (is risk-averse) if,
for every \omterm{def:b2:expected-utility:lottery}{lottery}, she weakly prefers its expected value for certain to the \omterm{def:b2:expected-utility:lottery}{lottery} itself:
$u(\mathbb{E}[\tilde z]) \ge \mathbb{E}[u(\tilde z)]$. She is \emph{risk-neutral} if this always holds
with equality, and \emph{risk-loving} if the inequality always goes the other way.
\end{definition}

\begin{proposition}[Risk aversion is concavity]\label{prop:b2:expected-utility:concave}
A person is risk-averse if and only if her utility of money is concave; risk-neutral if and only if it
is affine.
\end{proposition}

\begin{proof}
If $u$ is concave, Jensen's inequality gives $\mathbb{E}[u(\tilde z)] \le u(\mathbb{E}[\tilde z])$
for every distribution. Conversely, for two amounts $z, z'$ and $\theta \in [0,1]$, the \omterm{def:b2:expected-utility:lottery}{lottery} giving
$z$ with probability $\theta$ and $z'$ otherwise has expected value $\theta z + (1-\theta) z'$, so \omterm{def:b2:expected-utility:risk-aversion}{risk
aversion} says $u(\theta z + (1-\theta) z') \ge \theta u(z) + (1-\theta) u(z')$: this is the definition
of concavity. Affine: apply both inequalities.
\end{proof}

\begin{definition}[Certainty equivalent]\label{def:b2:expected-utility:certainty-equivalent}
The \emph{certainty equivalent}\index{certainty equivalent} $c(\tilde z)$ of a random amount is the
sure amount with the same expected utility: $u(c(\tilde z)) = \mathbb{E}[u(\tilde z)]$.
\end{definition}

It exists and is unique when $u$ is continuous and increasing (by the intermediate-value theorem, it
lies between the smallest and the largest prize), and it does not change when $u$ is replaced by
$\kappa_0 + \kappa_1 u$: it is measured in money, not in utils.

\begin{definition}[Risk premium]\label{def:b2:expected-utility:risk-premium}
The \emph{risk premium}\index{risk premium} of a random amount is $\rho(\tilde z) = \mathbb{E}[\tilde z]
- c(\tilde z)$: the amount of expected value the person would give up to replace the \omterm{def:b2:expected-utility:lottery}{lottery} by a sure
payment.
\end{definition}

\begin{proposition}[Premium and attitude]\label{prop:b2:expected-utility:premium-sign}
A person is risk-averse if and only if every \omterm{def:b2:expected-utility:risk-premium}{risk premium} she has is non-negative, and risk-neutral if
and only if every one is zero.
\end{proposition}

\begin{proof}
$u$ is increasing, so $c(\tilde z) \le \mathbb{E}[\tilde z] \iff u(c(\tilde z)) \le
u(\mathbb{E}[\tilde z]) \iff \mathbb{E}[u(\tilde z)] \le u(\mathbb{E}[\tilde z])$.
\end{proof}

\begin{example}[A square-root utility]\label{ex:b2:expected-utility:sqrt}
Let $u(z) = \sqrt z$, and let $\tilde z$ be \num{36} or \num{196} with probability $\tfrac12$ each. Its
expected value is $116$; its expected utility is $\tfrac12(6 + 14) = 10$; its \omterm{def:b2:expected-utility:certainty-equivalent}{certainty equivalent} is
$c = 10^2 = 100$; its \omterm{def:b2:expected-utility:risk-premium}{risk premium} is $116 - 100 = 16$.
The figure below draws it: the expected utility sits on the chord, below the curve,
and the \omterm{def:b2:expected-utility:certainty-equivalent}{certainty equivalent} is read off where the curve reaches that height.
\end{example}

\begin{omfigure}
\begin{tikzpicture}[econ]
\begin{axis}[omrecord, width=0.72\linewidth, height=0.46\linewidth, xmin=0, xmax=225, ymin=0,
    ymax=16, xlabel={Amount $z$}, ylabel={$u(z) = \sqrt{z}$}, grid=none,
    xtick={36,100,116,196}, ytick={6,10,14}, axis lines=left]
  \addplot[curve, domain=0:225, samples=120] {sqrt(x)};
  \draw[omThm, thick] (axis cs:36,6) -- (axis cs:196,14);
  \draw[guide] (axis cs:36,0) -- (axis cs:36,6);
  \draw[guide] (axis cs:196,0) -- (axis cs:196,14);
  \draw[guide] (axis cs:116,0) -- (axis cs:116,{sqrt(116)});
  \draw[guide] (axis cs:0,10) -- (axis cs:116,10);
  \draw[guide] (axis cs:100,0) -- (axis cs:100,10);
  \node[eqpoint] at (axis cs:116,10) {};
  \node[eqpoint] at (axis cs:100,10) {};
  \node[eqpoint] at (axis cs:116,{sqrt(116)}) {};
  \draw[omDef, very thick, <->] (axis cs:100,2.2) -- node[above, font=\footnotesize]{$\rho = 16$}
    (axis cs:116,2.2);
  \node[font=\footnotesize, anchor=north west] at (axis cs:118,9.8) {$\mathbb{E}[u] = 10$};
  \node[font=\footnotesize, anchor=south east] at (axis cs:114,{sqrt(116)+0.2})
    {$u(\mathbb{E}[\tilde z])$};
\end{axis}
\end{tikzpicture}

{\small A concave utility of money. The \omterm{def:b2:expected-utility:lottery}{lottery} pays \num{36} or \num{196}; its expected utility, $10$,
lies on the chord above the mean $116$. The curve reaches $10$ at $z = 100$, the \omterm{def:b2:expected-utility:certainty-equivalent}{certainty
equivalent}, so the \omterm{def:b2:expected-utility:risk-premium}{risk premium} is $116 - 100 = 16$.}
\end{omfigure}

Some lotteries are ranked by every person who prefers more money to less, whatever her attitude to
risk.

\begin{definition}[First-order stochastic dominance]\label{def:b2:expected-utility:fosd}
Let $F_\ell(z) = \sum_{i : z_i \le z} \ell_i$ be the distribution function of a \omterm{def:b2:expected-utility:lottery}{lottery} $\ell$ over
amounts of money. $\ell$ shows \emph{first-order stochastic dominance}\index{first-order stochastic
dominance} over $\ell'$ (it first-order dominates $\ell'$) if $F_\ell(z) \le F_{\ell'}(z)$ for every
$z$: for every amount, $\ell$ is at most as likely as $\ell'$ to pay that amount or less.
\end{definition}

\begin{proposition}[Dominance and increasing utility]\label{prop:b2:expected-utility:fosd}
$\ell$ first-order dominates $\ell'$ if and only if $\sum_i \ell_i u(z_i) \ge \sum_i \ell'_i u(z_i)$
for every increasing function $u$.
\end{proposition}

\begin{proof}
Put both lotteries on the same prizes $z_1 < \dots < z_n$ and write $F, F'$ for their distribution
functions. Summing by parts, with $F(z_n) = F'(z_n) = 1$,
\[
\sum_i (\ell_i - \ell'_i)\,u(z_i) = \sum_{i=1}^{n-1} \bigl(F'(z_i) - F(z_i)\bigr)\bigl(u(z_{i+1}) -
u(z_i)\bigr).
\]
If $F \le F'$ and $u$ is increasing, every term is $\ge 0$. Conversely, the increasing function $u(z)
= 1$ for $z > z_k$ and $0$ otherwise gives $\sum_i (\ell_i - \ell'_i)u(z_i) = F'(z_k) - F(z_k)$, which
must then be $\ge 0$ for each $k$.
\end{proof}

In the Allais pairs no option dominates the other: $\ell^{A}$ is more likely than $\ell^{B}$ to pay
$0$ ($F_A(0) = 0.01 > 0 = F_B(0)$) but less likely to pay \num{2400} or less ($0.67 < 1$). The theory
of risk attitudes is needed precisely where dominance is silent.

\section{Measuring risk aversion}\label{sec:b2:expected-utility:arrow-pratt}

Concavity says whether a person is risk-averse; it does not say how much. The curvature $u''$ is not
a measure, since it changes when $u$ is multiplied by a constant. Dividing by $u'$ removes the scale.

\begin{definition}[Coefficient of absolute risk aversion]\label{def:b2:expected-utility:ara}
For a twice differentiable increasing $u$, the \emph{coefficient of absolute risk
aversion}\index{coefficient of absolute risk aversion} at wealth $z$ is $\mathcal{A}(z) =
-u''(z)/u'(z)$, measured per unit of money.
\end{definition}

\begin{definition}[Coefficient of relative risk aversion]\label{def:b2:expected-utility:rra}
The \emph{coefficient of relative risk aversion}\index{coefficient of relative risk aversion} is
$\mathcal{R}(z) = z\,\mathcal{A}(z) = -z\,u''(z)/u'(z)$, a pure number.
\end{definition}

Both are unchanged by $u \mapsto \kappa_0 + \kappa_1 u$, as a measure built on the
\cref{prop:b2:expected-utility:affine} must be. The next result says what they measure.

\begin{proposition}[The small-risk approximation]\label{prop:b2:expected-utility:small-risk}
Let a person with wealth $z$ face a small risk $\tilde\varepsilon$ with mean $0$ and variance
$\sigma^2$. To second order in the size of the risk, her \omterm{def:b2:expected-utility:risk-premium}{risk premium} is
\[
\rho \approx \tfrac12\,\mathcal{A}(z)\,\sigma^2 .
\]
For a risk proportional to wealth, $z(1 + \tilde\varepsilon)$, the premium as a share of wealth is
$\rho/z \approx \tfrac12\,\mathcal{R}(z)\,\sigma^2$.
\end{proposition}

\begin{proof}
The premium solves $u(z - \rho) = \mathbb{E}[u(z + \tilde\varepsilon)]$. On the left, to first order
in $\rho$, $u(z - \rho) \approx u(z) - \rho\,u'(z)$. On the right, to second order,
$\mathbb{E}[u(z + \tilde\varepsilon)] \approx u(z) + \mathbb{E}[\tilde\varepsilon]\,u'(z) +
\tfrac12\mathbb{E}[\tilde\varepsilon^2]\,u''(z) = u(z) + \tfrac12\sigma^2 u''(z)$. Equating gives $\rho
\approx -\tfrac12\sigma^2 u''(z)/u'(z)$. (The premium is of the order of $\sigma^2$, so the term
$\rho^2 u''$ dropped on the left is of higher order.) For the proportional risk the variance is
$z^2\sigma^2$, so $\rho \approx \tfrac12\mathcal{A}(z) z^2\sigma^2$ and $\rho/z \approx
\tfrac12\mathcal{R}(z)\sigma^2$.
\end{proof}

The premium is proportional to the variance, not to the standard deviation. A person with a twice
differentiable utility is therefore close to risk-neutral towards very small risks: halving a risk
divides its premium by four. This is called second-order \omterm{def:b2:expected-utility:risk-aversion}{risk aversion}, and it returns in the last
section.

Two families of utility functions hold one of the coefficients fixed.

\begin{definition}[Constant absolute risk aversion]\label{def:b2:expected-utility:cara}
A utility of money shows \emph{constant absolute risk aversion}\index{constant absolute risk aversion}
(CARA) with coefficient $a > 0$ if $u(z) = -e^{-az}/a$; then $\mathcal{A}(z) = a$ for every $z$.
\end{definition}

\begin{definition}[Constant relative risk aversion]\label{def:b2:expected-utility:crra}
A utility of money shows \emph{constant relative risk aversion}\index{constant relative risk aversion}
(CRRA) with coefficient $\gamma > 0$ if
\[
u(z) = \frac{z^{1-\gamma}}{1-\gamma} \quad (\gamma \ne 1), \qquad u(z) = \ln z \quad (\gamma = 1),
\]
for $z > 0$; then $\mathcal{R}(z) = \gamma$ and $\mathcal{A}(z) = \gamma / z$.
\end{definition}

The logarithm is the limit case: $(z^{1-\gamma} - 1)/(1-\gamma)$, which differs from the CRRA utility
only by a constant and a factor, tends to $\ln z$ as $\gamma \to 1$. For $\gamma > 1$ the utility is
negative and bounded above by $0$, a fact that matters below.

\begin{omfigure}
\begin{tikzpicture}[econ]
\begin{axis}[omrecord, width=0.7\linewidth, height=0.46\linewidth, xmin=0.2, xmax=3, ymin=-3,
    ymax=2.2, xlabel={Wealth $z$ (in units of a reference wealth)}, ylabel={$u(z)$}, grid=none,
    legend pos=south east]
  \addplot[omDef, thick, domain=0.2:3, samples=80] {(x^0.5 - 1)/0.5};
  \addplot[omThm, thick, domain=0.2:3, samples=80] {ln(x)};
  \addplot[omMeth, thick, domain=0.2:3, samples=80] {(x^(-1) - 1)/(-1)};
  \addplot[omProp, thick, domain=0.2:3, samples=80] {(x^(-4) - 1)/(-4)};
  \legend{$\gamma = 0.5$, $\gamma = 1$, $\gamma = 2$, $\gamma = 5$}
\end{axis}
\end{tikzpicture}

{\small \omterm{def:b2:expected-utility:crra}{Constant relative risk aversion}, $u(z) = (z^{1-\gamma} - 1)/(1-\gamma)$, scaled so that every
curve passes through $(1, 0)$ with slope $1$. The higher $\gamma$, the faster marginal utility falls
and the flatter the curve above the reference wealth.}
\end{omfigure}

\begin{proposition}[What stays constant]\label{prop:b2:expected-utility:cara-crra}
\begin{enumerate}
  \item Under CARA, the \omterm{def:b2:expected-utility:risk-premium}{risk premium} of an added risk $\tilde\varepsilon$ does not depend on wealth:
        $c(z + \tilde\varepsilon) = z + c(\tilde\varepsilon)$, exactly.
  \item Under CRRA, the \omterm{def:b2:expected-utility:risk-premium}{risk premium} of a proportional risk is proportional to wealth: $c(z(1 +
        \tilde\varepsilon)) = z\,c(1 + \tilde\varepsilon)$, exactly.
\end{enumerate}
\end{proposition}

\begin{proof}
(1) $u(z + x) = -e^{-az}e^{-ax}/a = e^{-az}u(x)$: adding $z$ multiplies every utility by the same
positive number, which does not change the ranking, so the sure amount matching $z + \tilde\varepsilon$
is $z$ plus the one matching $\tilde\varepsilon$. (2) For $\gamma \ne 1$, $u(zx) = z^{1-\gamma}u(x)$;
for $\gamma = 1$, $\ln(zx) = \ln z + \ln x$. Either way scaling wealth by $z$ is an increasing affine
change of $u$ and leaves the ranking unchanged.
\end{proof}

\begin{example}[The premium as wealth rises]\label{ex:b2:expected-utility:wealth}
Let a person face a 50--50 gain or loss of \usd{5000}. With CRRA $\gamma = 2$, $u(z) = -1/z$ and
\[
\mathbb{E}[u] = -\tfrac12\Bigl(\frac{1}{z - h} + \frac{1}{z + h}\Bigr) = -\frac{z}{z^2 - h^2}, \qquad
c = \frac{z^2 - h^2}{z}, \qquad \rho = \frac{h^2}{z},
\]
with $h = \num{5000}$: \usd{1250} at a wealth of \usd{20000}, \usd{500} at \usd{50000}. With CARA and $a =
\num{0.0001}$ per dollar (the absolute \omterm{def:b2:expected-utility:risk-aversion}{risk aversion} of CRRA $\gamma = 2$ at \usd{20000}), $\rho = \ln\cosh(ah)/a = \ln\cosh(0.5)/\num{0.0001} \approx \usd{1201}$ at every wealth
(\cref{exo:b2:expected-utility:8} derives the formula).
The figure below plots both. Under CARA a household with \usd{1000000} and one
with \usd{15000} pay the same premium to shed the same risk; part IV of the weekend problem turns this
difference between the two families into a test.
\end{example}

\begin{omfigure}
\begin{tikzpicture}[econ]
\begin{axis}[omrecord, width=0.7\linewidth, height=0.44\linewidth, xmin=10000, xmax=60000, ymin=0,
    ymax=2600, xlabel={Wealth $z$ (dollars)}, ylabel={Risk premium $\rho$ (dollars)}, grid=none,
    legend pos=north east, xtick={10000,20000,30000,40000,50000,60000}]
  \addplot[omDef, thick, domain=10000:60000, samples=100] {25e6/x};
  \addplot[omThm, thick, dashed, domain=10000:60000, samples=2] {ln(cosh(0.5))/0.0001};
  \legend{CRRA $\gamma = 2$, CARA $a = \num{0.0001}$}
  \draw[guide] (axis cs:20000,0) -- (axis cs:20000,1250);
\end{axis}
\end{tikzpicture}

{\small The \omterm{def:b2:expected-utility:risk-premium}{risk premium} of a 50--50 gain or loss of \usd{5000}. Under CRRA with $\gamma = 2$ it is
$h^2/z$ and falls as wealth rises; under CARA it is the same at every wealth. The two utilities have
the same absolute \omterm{def:b2:expected-utility:risk-aversion}{risk aversion} at \usd{20000} (dotted line), where their premiums nearly agree.}
\end{omfigure}

The coefficients also compare two people.

\begin{proposition}[Comparing risk aversion (Pratt)]\label{prop:b2:expected-utility:pratt}
Let $u$ and $v$ be twice differentiable, increasing utilities of money. The following are equivalent:
\begin{enumerate}
  \item $\mathcal{A}_u(z) \ge \mathcal{A}_v(z)$ for every $z$;
  \item $u = \varphi \circ v$ for an increasing concave function $\varphi$;
  \item for every wealth and every risk, the \omterm{def:b2:expected-utility:risk-premium}{risk premium} under $u$ is at least the one under $v$.
\end{enumerate}
\end{proposition}

\begin{proof}
(1)$\iff$(2). Since $v$ is increasing, $\varphi = u \circ v^{-1}$ exists, is increasing and twice
differentiable, and $u = \varphi(v)$. Differentiating twice, $u' = \varphi'(v)v'$ and $u'' =
\varphi''(v)v'^2 + \varphi'(v)v''$, so
\[
\mathcal{A}_u = -\frac{u''}{u'} = \mathcal{A}_v - \frac{\varphi''(v)}{\varphi'(v)}\,v' ,
\]
and $\mathcal{A}_u \ge \mathcal{A}_v$ everywhere if and only if $\varphi'' \le 0$. (2)$\Rightarrow$(3).
By Jensen, $u(c_u) = \mathbb{E}[\varphi(v(\tilde z))] \le \varphi(\mathbb{E}[v(\tilde z)]) =
\varphi(v(c_v)) = u(c_v)$, so $c_u \le c_v$ and $\rho_u \ge \rho_v$. (3)$\Rightarrow$(1). Apply (3) to
small risks: by \cref{prop:b2:expected-utility:small-risk}, $\tfrac12\mathcal{A}_u(z)\sigma^2 \ge
\tfrac12\mathcal{A}_v(z)\sigma^2$ up to terms of higher order, which gives (1) as $\sigma \to 0$.
\end{proof}

\section{Where the model is tested, and where it breaks}\label{sec:b2:expected-utility:tests}

\begin{model}{Expected utility}
\assumes A person ranks lotteries by a complete, transitive and continuous preference that satisfies
independence, and cares only about the reduced \omterm{def:b2:expected-utility:lottery}{lottery} over final outcomes.
\predicts Choices are linear in probabilities: straight parallel indifference lines in the triangle;
the same choice in both questions of an Allais pair; \omterm{def:b2:expected-utility:risk-premium}{risk premiums} that shrink like the variance for
small risks; full insurance when it is sold at its expected cost.
\tested In laboratory choices between gambles, where most respondents break it in the Allais pairs
(Kahneman and Tversky 1979); in home and flood insurance, where deductible choices imply implausibly
high \omterm{def:b2:expected-utility:risk-aversion}{risk aversion} (Sydnor 2010; Collier et al.\ 2018).
\breaks Near certainty (the certainty effect of the Allais pairs), and for small stakes, where the
curvature needed to explain the choices implies absurd choices over larger stakes.
\end{model}

\subsection{The Allais pairs}

\begin{proposition}[Common consequence]\label{prop:b2:expected-utility:common-consequence}
Under expected utility, $\ell^{A} \succ \ell^{B}$ if and only if $\ell^{C} \succ \ell^{D}$ in the
lotteries of \cref{ex:b2:expected-utility:allais-lotteries}.
\end{proposition}

\begin{proof}
Directly: $U(\ell^{A}) - U(\ell^{B}) = 0.01\,u(0) + 0.66\,u(\num{2400}) + 0.33\,u(\num{2500}) -
u(\num{2400}) = 0.01\,u(0) - 0.34\,u(\num{2400}) + 0.33\,u(\num{2500})$, and $U(\ell^{C}) - U(\ell^{D})
= 0.67\,u(0) + 0.33\,u(\num{2500}) - 0.66\,u(0) - 0.34\,u(\num{2400})$ is the same number. Or by
independence: both pairs mix $\ell^{*}$ against $\delta_2$ with weight $0.34$.
\end{proof}

\begin{proposition}[Common ratio]\label{prop:b2:expected-utility:common-ratio}
Let $\ell^{A'}$ pay \num{4000} with probability $0.8$, $\ell^{B'}$ pay \num{3000} for sure,
$\ell^{C'}$ pay \num{4000} with probability $0.2$ and $\ell^{D'}$ pay \num{3000} with probability
$0.25$ (otherwise nothing). Under expected utility, $\ell^{A'} \succ \ell^{B'} \iff \ell^{C'} \succ
\ell^{D'}$.
\end{proposition}

\begin{proof}
$\ell^{C'} = 0.25\,\ell^{A'} + 0.75\,\delta_0$ and $\ell^{D'} = 0.25\,\ell^{B'} + 0.75\,\delta_0$, where
$\delta_0$ pays nothing; apply independence with $\theta = 0.25$.
\end{proof}

\begin{record}[The Allais pairs]
The pairs are variants of a counter-example put forward by Allais (1953). Kahneman and Tversky (1979)
asked them with moderate prizes. In the common-consequence pair (Problems 1 and 2, 72 respondents),
82 per cent chose $\ell^{B}$, the sure \num{2400}, and 83 per cent chose $\ell^{C}$.
% ledger: b2:kt1979-p1, b2:kt1979-p2
In the common-ratio pair (Problems 3 and 4, 95 respondents), 80 per cent chose the sure \num{3000} and
65 per cent chose the \num{4000} at probability $0.2$.
% ledger: b2:kt1979-p3, b2:kt1979-p4
Both modal patterns contradict \cref{prop:b2:expected-utility:common-consequence} and
\cref{prop:b2:expected-utility:common-ratio}. What the record says: people value the step from a
probability of $0.99$ to certainty far more than the step from $0.33$ to $0.34$; the theory, linear in
probabilities, gives the two steps of $0.01$ the same weight. Kahneman and Tversky named this the
certainty effect. The choices were hypothetical, a limit the authors discuss: the method assumes that people know how
they would choose.
\end{record}

In the triangle, the modal choices say that the indifference line through $\ell^{B}$ is steeper than
the segment from $\ell^{B}$ to $\ell^{A}$ (slope $33$), while the line through $\ell^{D}$ is flatter
than the parallel segment from $\ell^{D}$ to $\ell^{C}$. Indifference curves that are steeper on the
left of the triangle than on its right are said to \emph{fan out}; \cref{exo:b2:expected-utility:11}
works the geometry. Chapter~\ref{ch:b2:behavioural-evidence} takes up the models built to fit this
pattern and their own tests.

\subsection{Small stakes}

Turn down a coin toss that loses \usd{100} or wins \usd{110}? Many people do, and it seems a mild
caution. Expected utility turns it into a strong statement.

\begin{proposition}[A calibration bound]\label{prop:b2:expected-utility:calibration}
Let $u$ be increasing, concave and differentiable. If a person turns down the 50--50 gamble that loses
\usd{100} or gains \usd{110} at every level of wealth, then for every $z$,
\[
u'(z + 210) \le \tfrac{10}{11}\,u'(z), \qquad\text{hence}\qquad u'(z + 210k) \le
\bigl(\tfrac{10}{11}\bigr)^{k} u'(z) \quad (k = 1, 2, \dots),
\]
and she turns down every 50--50 gamble that loses \usd{2400}, however large the gain.
\end{proposition}

\begin{proof}
Turning the gamble down at wealth $x$ means $\tfrac12 u(x - 100) + \tfrac12 u(x + 110) \le u(x)$,
that is, $u(x + 110) - u(x) \le u(x) - u(x - 100)$. By concavity, the slope of a chord lies between the
derivatives at its ends: $u(x + 110) - u(x) \ge 110\,u'(x + 110)$ and $u(x) - u(x - 100) \le
100\,u'(x - 100)$. So $110\,u'(x + 110) \le 100\,u'(x - 100)$; put $x = z + 100$. Iterate for the
second inequality. For the gain $G$, cut $[z, z + G]$ into pieces of length $210$; on the $k$-th piece
($k = 0, 1, \dots$) the slope is at most $u'(z + 210k)$, so
\[
u(z + G) - u(z) \le \sum_{k \ge 0} 210\,\bigl(\tfrac{10}{11}\bigr)^{k} u'(z) = 210 \times 11\,u'(z) =
\num{2310}\,u'(z).
\]
For the loss, concavity gives $u(z) - u(z - \num{2400}) \ge \num{2400}\,u'(z)$. Then $\tfrac12 u(z -
\num{2400}) + \tfrac12 u(z + G) - u(z) \le \tfrac12(\num{2310} - \num{2400})\,u'(z) < 0$.
\end{proof}

The bound says that every \usd{210} of extra wealth cuts the value of a marginal dollar by at least a
tenth. Over \usd{21000} ($k = 100$) the factor is $(10/11)^{100} \approx \num{7.3e-5}$: a dollar at
% code: exercises/solutions.py ex_calibration
the higher wealth is worth less than one thirteen-thousandth of a dollar at the lower one. No one's
marginal utility of money behaves like that over a range of \usd{21000}. The calibration result of
Rabin (2000) sharpens the bound in the same way; its conclusion is that, under expected utility,
aversion to small gambles cannot come from the curvature of the utility of wealth.
% ledger: b2:rabin2000

\begin{example}[The curvature it takes]\label{ex:b2:expected-utility:rabin-gamma}
With CRRA utility and a wealth of \usd{100000}, the coefficient that makes the \usd{100}/\usd{110}
gamble a tie is $\gamma \approx 91$ (the small-risk approximation, mean gain $5$ against $\tfrac12
\tfrac{\gamma}{z}\,105^2$, gives the same $91$). Such a person turns down a 50--50 gamble that loses
% code: exercises/solutions.py ex_rabin_gamma
\usd{769} whatever it pays: with $\gamma > 1$, utility is bounded above by $0$, so she rejects it as
soon as $\tfrac12 u(z - L) < u(z)$, that is, $L > z\bigl(1 - 2^{-1/(\gamma - 1)}\bigr)$. With
$\gamma = 2$ instead, the premium of the coin toss at the same wealth is about $\tfrac12 \times
(2/\num{100000}) \times 105^2 \approx \usd{0.11}$, far below its mean gain of \usd{5}: she takes it.
\end{example}

\begin{measure}[Reading risk aversion off insurance deductibles]
A deductible is the part of a claim the household pays itself. A lower deductible costs a higher
premium. Suppose a household pays an extra premium $\nu$ to lower its deductible from $d_h$ to
$d_\ell$, and it files a claim in a year with probability $\mu$. The expected saving is $\mu(d_h -
d_\ell)$, and the ratio $\nu / (\mu(d_h - d_\ell))$ is the load: what each dollar of expected
saving costs. Under CARA, the household's choice brackets $a$: it chose the low deductible only if
\[
e^{a\,\nu}\bigl(1 - \mu + \mu e^{a d_\ell}\bigr) \le 1 - \mu + \mu e^{a d_h} ,
\]
which sets a lower bound on $a$, and a \omterm{def:b2:expected-utility:rra}{coefficient of relative risk aversion} follows once a wealth is
chosen. The claim rate comes from the insurer's records for similar homes. Sydnor (2010), with the
records of a large US home insurer, found that most customers choose low deductibles at premiums far
above their expected value, which a standard model of \omterm{def:b2:expected-utility:risk-aversion}{risk aversion} fits only with implausibly large
coefficients.
% ledger: b2:sydnor2010
Collier et al.\ (2018) studied the flood-insurance choices of over \num{100000} US households: over 94
% ledger: b2:cskm2018-households, b2:cskm2018-lowdeduct
per cent chose one of the two lowest of six deductibles, the median premium was 2.7 times the expected
% ledger: b2:cskm2018-load
claim payment, and the deductible choices imply a median \omterm{def:b2:expected-utility:rra}{coefficient of relative risk aversion} of 108.
% ledger: b2:cskm2018-rra
\end{measure}

The two kinds of evidence point the same way. A coefficient near $100$ explains a low deductible or a
refused coin toss, and predicts absurd refusals of larger gambles; a coefficient of $2$ predicts that
the coin toss is taken and that a low deductible at a load of $2.7$ is not bought. One utility of final
wealth cannot do both. The alternatives that fit small stakes (losses weighed more than gains, measured from a
reference point, and probabilities weighted non-linearly) are taken up with their evidence in
\cref{ch:b2:behavioural-evidence}.

\begin{remark}[What is kept]\label{rem:b2:expected-utility:kept}
For decisions whose stakes are a sizeable share of wealth (insurance of a home or a harvest, saving
for retirement, the choice of a portfolio) expected utility with a moderate coefficient remains the
model of the next three chapters, because it gives sharp predictions that the data can check: full
insurance at a fair price, a premium that shrinks with wealth under decreasing absolute \omterm{def:b2:expected-utility:risk-aversion}{risk aversion},
a portfolio share that rises with the excess return. Where its predictions fail, the failure is
recorded with the model.
\end{remark}

\section{Exercises}\label{sec:b2:expected-utility:exercises}

\begin{exercise}[$\star$]\label{exo:b2:expected-utility:1}
Read the opening chart. In each pair, a respondent who obeys expected utility chooses either the less
risky option twice or the riskier option twice. What is the smallest percentage of respondents who
broke expected utility in the first pair? In the second? Why is it only a lower bound?
\end{exercise}

\begin{exercise}[$\star$]\label{exo:b2:expected-utility:2}
A \omterm{def:b2:expected-utility:compound-lottery}{compound lottery} gives, with probability $\tfrac12$ each, the \omterm{def:b2:expected-utility:lottery}{lottery} $\ell$ (0 with probability $0.2$,
\num{100} with probability $0.8$) or the \omterm{def:b2:expected-utility:lottery}{lottery} $\ell'$ (0 with probability $0.6$, \num{100} with
probability $0.4$). Find the reduced \omterm{def:b2:expected-utility:lottery}{lottery} and its expected value. With $u(z) = \sqrt z$, find its
expected utility and its \omterm{def:b2:expected-utility:certainty-equivalent}{certainty equivalent}.
\end{exercise}

\begin{exercise}[$\star$]\label{exo:b2:expected-utility:3}
With $u(z) = \ln z$, a \omterm{def:b2:expected-utility:lottery}{lottery} pays \num{50} or \num{200} with probability $\tfrac12$ each. Find its
expected value, \omterm{def:b2:expected-utility:certainty-equivalent}{certainty equivalent} and \omterm{def:b2:expected-utility:risk-premium}{risk premium}. Compare the premium with the small-risk
approximation at the mean.
\end{exercise}

\begin{exercise}[$\star$]\label{exo:b2:expected-utility:4}
Compute $\mathcal{A}(z)$ and $\mathcal{R}(z)$ for $u(z) = -e^{-0.002z}$, $\sqrt z$, $\ln z$ and $-1/z$,
and their values at $z = \num{1000}$. Which are CARA, which CRRA? Which shows absolute \omterm{def:b2:expected-utility:risk-aversion}{risk aversion}
that falls with wealth?
\end{exercise}

\begin{exercise}[$\star\star$]\label{exo:b2:expected-utility:5}
In Problems 1 and 2, normalise $u(0) = 0$ and $u(\num{2400}) = 1$ (why is this allowed?). Find the set
of values of $u(\num{2500})$ for which an expected-utility maximiser chooses $\ell^{B}$ in Problem 1,
and the set for which she chooses $\ell^{C}$ in Problem 2. Do the same for Problems 3 and 4 with $u(0)
= 0$, $u(\num{4000}) = 1$. Which pair of observed choices would falsify expected utility for a single
respondent?
\end{exercise}

\begin{exercise}[$\star\star$]\label{exo:b2:expected-utility:6}
The \omterm{def:b2:expected-utility:lottery}{lottery} $\ell$ pays 0, \num{50} or \num{100} with probabilities $0.2$, $0.3$, $0.5$; $\ell'$ pays
them with probabilities $0.3$, $0.3$, $0.4$; $\ell''$ with probabilities $0.1$, $0.6$, $0.3$. Which of
these lotteries first-order dominates which? For each pair not ranked, find one increasing utility
that prefers each \omterm{def:b2:expected-utility:lottery}{lottery} (try $u(z) = z$ and $u(z) = \sqrt z$).
\end{exercise}

\begin{exercise}[$\star\star$]\label{exo:b2:expected-utility:7}
A household with CRRA utility, $\gamma = 3$, and wealth \usd{50000} faces a 50--50 gain or loss of
\usd{5000}. Predict its \omterm{def:b2:expected-utility:risk-premium}{risk premium} with the small-risk approximation, then compute it exactly. By how
much, roughly, does the premium change if its wealth doubles, and why?
\end{exercise}

\begin{exercise}[$\star\star$]\label{exo:b2:expected-utility:8}
Show that under CARA with coefficient $a$ the \omterm{def:b2:expected-utility:risk-premium}{risk premium} of a 50--50 gain or loss of $h$ is
$\ln\cosh(ah)/a$ at every wealth. Evaluate it for $a = \num{0.0005}$ per dollar and $h = \usd{1000}$,
and compare with the small-risk approximation.
\end{exercise}

\begin{exercise}[$\star\star$]\label{exo:b2:expected-utility:9}
A household pays \usd{54} a year to lower its home-insurance deductible from \usd{2000} to \usd{1000};
it files a claim in a year with probability $0.02$. Find the expected saving and the load, and compare
the load with the median in the measure box. Using the inequality of the box, find the smallest CARA
coefficient that explains the choice, and the \omterm{def:b2:expected-utility:rra}{coefficient of relative risk aversion} at a wealth of
\usd{200000}. Is the order of magnitude that of the box?
\end{exercise}

\begin{exercise}[$\star\star\star$]\label{exo:b2:expected-utility:10}
Generalise \cref{prop:b2:expected-utility:calibration}: a person with an increasing, concave and
differentiable $u$ turns down, at every wealth, the 50--50 gamble that loses $l$ or gains $g$, with $0
< l < g$. Prove that $u'(z + l + g) \le (l/g)\,u'(z)$, and that she turns down every 50--50 gamble that
loses more than $(l + g)\,g/(g - l)$, whatever it gains. Check the bound \num{2310} for $l = 100$, $g =
110$, and say what observation would show that a person's aversion to the small gamble does not come
from the curvature of a utility of wealth.
\end{exercise}

\begin{exercise}[$\star\star\star$]\label{exo:b2:expected-utility:11}
In the probability triangle with prizes $z_1 < z_2 < z_3$, show that a person with utility $u = \varphi
\circ v$, $\varphi$ increasing and concave, has indifference lines at least as steep as those of a
person with utility $v$. Show that the modal choices of Problems 1 and 2 require an indifference line
through $\ell^{B}$ steeper than $33$ and one through $\ell^{D}$ flatter than $33$. What pattern of
choices, over many pairs of lotteries placed across the triangle, would show the lines fanning out
rather than parallel?
\end{exercise}

\begin{exercise}[$\star\star\star$]\label{exo:b2:expected-utility:12}
In a list of ten rows, a respondent chooses in each row between a safer \omterm{def:b2:expected-utility:lottery}{lottery} (\dal{20} or
\dal{16}) and a riskier one (\dal{38} or \dal{1}); in row $k$ the better prize of each comes with
probability $k/10$. Prove that any expected-utility maximiser with an increasing utility switches at
most once, from the safer to the riskier \omterm{def:b2:expected-utility:lottery}{lottery}, as $k$ rises. A respondent takes the safer \omterm{def:b2:expected-utility:lottery}{lottery}
in rows 1 to 6 and the riskier in rows 7 to 10: with CRRA utility of the prize, which interval of
$\gamma$ is consistent with her choices (solve numerically)? What would a respondent who switches back
and forth tell you?
\end{exercise}

\section{Problem: A Brenland grain farmer's harvest}\label{sec:b2:expected-utility:problem}

\begin{problem}[{Weekend problem --- a farmer in Brenland, an invented economy, whose income depends
on the rain: the harvest as a lottery, a choice of seed, an offer of drought insurance, and what her
savings change}]
\label{pb:b2:expected-utility:1}
A farmer's income is \dal{60000} after a good harvest, which comes with probability $0.7$, and
\dal{20000} after a drought, with probability $0.3$. She has no other income and maximises the
expected utility $\mathbb{E}[\ln \tilde z]$ of her income $\tilde z$. Amounts are in dalens; round to
the nearest dalen.

\medskip
\noindent\textbf{Part I --- The harvest as a lottery.}
\begin{enumerate}
  \item Find the mean and the standard deviation of her income.
  \item Find her expected utility.
  \item Find the \omterm{def:b2:expected-utility:certainty-equivalent}{certainty equivalent} of the harvest.
  \item Find its \omterm{def:b2:expected-utility:risk-premium}{risk premium}.
  \item Compute the small-risk approximation of the premium at the mean. Why is it so far off?
  \item A neighbour with the same harvest has $u(z) = \sqrt z$. Find his \omterm{def:b2:expected-utility:certainty-equivalent}{certainty equivalent} and \omterm{def:b2:expected-utility:risk-premium}{risk
        premium}, and explain the ranking with \cref{prop:b2:expected-utility:pratt}.
\end{enumerate}

\medskip
\noindent\textbf{Part II --- Choosing a seed.}
The farmer's usual seed is the one above (call it T). A hardier seed S yields \dal{55000} in a good
year and \dal{30000} in a drought.
\begin{enumerate}[resume]
  \item Find the mean income with S. Which seed would a risk-neutral farmer choose?
  \item Does either seed first-order dominate the other?
  \item Find her expected utility and \omterm{def:b2:expected-utility:certainty-equivalent}{certainty equivalent} with S. Which seed does she choose?
  \item For which CRRA coefficient $\gamma$ is a farmer indifferent between S and T (solve
        numerically)? Which seed do more risk-averse farmers choose?
  \item A seed V yields \dal{61000} and \dal{20000}. Without computing, rank V and T for every farmer
        who prefers more income to less.
  \item She plants half her land with T and half with S. Find her income in each state, her expected
        utility and her \omterm{def:b2:expected-utility:certainty-equivalent}{certainty equivalent}. Compare with planting all S, and explain.
\end{enumerate}

\medskip
\noindent\textbf{Part III --- An offer of drought insurance.}
She plants T. An insurer offers to pay \dal{40000} in a drought, for a premium paid in every year.
\begin{enumerate}[resume]
  \item Find the premium equal to the expected payment (the fair premium) and her income in each state
        if she buys at that premium. Does she buy?
  \item Find the largest premium she is willing to pay for this cover.
  \item The loading is the premium divided by the expected payment, minus one. Find the largest loading
        she accepts.
  \item The insurer now sells any amount of cover $\chi$ (dalens paid in a drought) at a price $\psi$ per
        dalen of cover. Write her income in each state and show that her best cover is $\chi^{*} =
        (\num{18000} - \num{32000}\,\psi)/(\psi (1-\psi))$, as long as this lies between $0$ and \num{40000}.
  \item Find $\chi^{*}$ at the fair price $\psi  = 0.3$, and explain the result.
  \item Find $\chi^{*}$ at a loading of 20 per cent ($\psi  = 0.36$).
  \item Above which price does she buy no cover at all? Express it as a loading, and explain why it is
        higher than the answer to question 15.
\end{enumerate}

\medskip
\noindent\textbf{Part IV --- Savings, and a test.}
\begin{enumerate}[resume]
  \item A second farmer has CARA utility with the coefficient $a$ that gives her the same \omterm{def:b2:expected-utility:risk-premium}{risk premium}
        for the harvest as the first farmer. Find $a$ numerically, and compare it with $\mathcal{A}$ of
        $\ln z$ at the mean income.
  \item Both farmers now hold \dal{30000} of savings, paid in both states. Find the \omterm{def:b2:expected-utility:risk-premium}{risk premium} of the
        harvest for each.
  \item Find, for each, the largest premium and the largest loading they accept for the cover of Part
        III.
  \item An insurer offers the cover at a loading of 30 per cent to farmers with and without savings.
        What does each model predict about who buys? What data on take-up would reject each?
  \item State the result: the harvest's \omterm{def:b2:expected-utility:risk-premium}{risk premium}, and the largest loading the farmer accepts for
        full drought cover, without savings and with \dal{30000} of savings.
\end{enumerate}
\end{problem}
